\section{Conclusions}
\label{sec:conclusions}

We recall that the goal of the project was to design a function generator capable of producing output waveforms with frequencies up to \(1 MHz\). \\
Laboratory tests have shown that the proposed design fully meets the specified requirements. \\
Nevertheless, some parts of the system could be optimized, and several possible upgrades may be implemented to enhance the overall functionality of the design without significantly increasing the hardware complexity of the chip.

\subsection{Criticality}
Even though it has been verified with the timing analyzer that the dynamic clock generated by the frequency divider instantiated in the DDFS module does not cause any timing violations, this verification has been carried out by considering the clock source in phase with the main clock, without taking into account the \(t_{co}\) of the flip-flop. \\
To avoid potential issues, the divided clock can be considered asynchronous and the input can be protected using a flip-flop chain. \\
A more professional solution can be obtained by using a \texttt{clock\_enable} signal, which identifies, via a counter synchronized with the main clock, the \(n\)-th clock cycle and synchronizes the module using this signal, where \(n\) is controlled by the \texttt{freq\_control} input. \\
A possible Verilog implementation of the clock enable signal is shown in \autoref{lst:clk_ena}.

\vspace{0.5cm}
\begin{lstlisting}[language=Verilog, caption=Implementation of clock enable signal with frequency divider, 
label=lst:clk_ena ]
    integer cont;
    always @(posedge clk) begin
        if(rst_n == 0) begin
            //reset condition
            cont <= 0;
            clk_ena <= 1'b0;
        end
        else begin
            if(cont >= n) begin
                cont <= 0;
                clk_ena <= 1'b1;
            end
            else begin
                cont <= cont + 1;
                clk_ena <= 1'b0;
            end
        end
    end
\end{lstlisting}
\vspace{0.5cm}

\noindent
With the \texttt{clock\_ena} signal, all the sequential blocks synchronized with the divided frequency clock become like the block in \autoref{lst:clk_div_ena} 

\vspace{0.5cm}
\begin{lstlisting}[language=Verilog, caption=Avoid using divided clock, 
label=lst:clk_div_ena ]
    always @(posedge clk) begin
        if(clk_ena == 1) begin
            //divided clock operation
        end
        else begin
            //store the signals
        end
    end
\end{lstlisting}
\vspace{0.5cm}

\newpage

\subsection{Possible optimizations}
The first hardware optimization that can be done is inside the \texttt{DDFS\_frequency\_converter} where a large number of multipliers have been instantiated.

\vspace{0.5cm}
\begin{lstlisting}[language=Verilog, caption=Part of the calculations performed to obtain the frequency word value, 
label=lst:mult_inst_non_opt ]
    always @(*) begin
        
        if(final_freq <= (CLK_FREQ/(1000000*MIN_LUT_SAMPLE))) begin
            freq_control = 3'd6;
            fw = (fw_1000000 >= 0) ? fw_1000000 : 7'd0;
        end 
        else if(final_freq <= (CLK_FREQ/(100000*MIN_LUT_SAMPLE))) begin
            freq_control = 3'd5;
            fw = (fw_100000 >= 0) ? fw_100000 : 7'd0;
        end
        .
        .
        .
        
        else begin
            //default condition, max frequency
            freq_control = 3'd0;
            fw = 7'd0;
        end
        
    end

    always @(posedge clk) begin
        ...
    
        fw_1000000 <= (mult_1000000*K)>>54;
        fw_100000 <= (mult_100000*K)>>54;
        fw_10000 <= (mult_10000*K)>>54;
        fw_1000 <= (mult_1000*K)>>54;
        fw_100 <= (mult_100*K)>>54;
        fw_10 <= (mult_10*K)>>54;
        fw_2 <= (mult_2*K)>>54;
    end
\end{lstlisting}
\vspace{0.5cm}

\noindent
In \autoref{lst:mult_inst_non_opt}, is calculated the value of \texttt{f\_w} for every clock value that goes into the DDFS module. \\
\texttt{fw\_2} is the frequency word calculated with the clock divided by two, \texttt{fw\_10} is the same with the clock divided by ten and so on.

\newpage
\noindent
The problem is that \texttt{mult\_*} is a \(64 bit\) value, so the multipliers instantiated are big combinatorial network. \\
To reduce the number of multiplier instantiated the input of the operation can be multiplexed as shown in \autoref{lst:mult_inst_opt}

\vspace{0.5cm}
\begin{lstlisting}[language=Verilog, caption=Example to optimize the calculation of frequency word value, 
label=lst:mult_inst_opt ]
    always (*) begin
    
        if(final_freq <= (CLK_FREQ/(1000000*MIN_LUT_SAMPLE))) begin
            freq_control = 3'd6;
            mult = mult_1000000;
        end 
        else if(final_freq <= (CLK_FREQ/(100000*MIN_LUT_SAMPLE))) begin
            freq_control = 3'd5;
            mult = mult_100000;
        end
        .
        .
        .
        
        else begin
            //default condition, max frequency
            freq_control = 3'd0;
            mult = mult_2;
        end

    end

    always (posedge clk) begin
        ...
        
        fw <= (mult*K)>>54;
    end
\end{lstlisting}
\vspace{0.5cm}

\noindent
This optimization instantiates a multiplexer and can reduce the maximum frequency of the project; to overcome this problem, the multiplexing operation can be pipelined by modifying the combinatorial block into a synchronous block. \\
The pipeline operation increases the frequency but also the throughput of the system.\\
Another optimization that can be done is to allow the \texttt{ddfs} module to work with its own clock source and control it asynchronously using a FIFO memory.\\
This is because the \texttt{ddfs} module can work at a frequency much higher than the other modules, due to its hardware simplicity, and the other modules are used to interface with the operator, so they can produce the outputs even after some milliseconds.\\
This upgrade allows the frequency generator to produce higher-frequency waveforms without increasing the complexity of the circuit. \\
Another possible optimization concerns the LUTs that store the contents of the screens displayed on the LCD. For future upgrades, some spaces are left blank to allow the addition of more pages. However, in the final version, all unused blank spaces can be removed, reducing the amount of instantiated LUT memory.

\subsection{Possible upgrades}
A simple upgrade that can be made is to instantiate two additional \texttt{ddfs} modules within the project and assign each module a dedicated button on the keypad. \\
By doing this, we can use all three DAC outputs to generate dedicated waveforms. \\
A more advanced update could also allow the modules to communicate with each other, enabling their outputs to be synchronized, either in phase or with a programmed phase shift. \\
Furthermore, it is possible to implement a method, via USB communication or SD card, to upload samples of customizable waveforms for display. These samples are then saved into the board's internal FLASH memory and can be displayed when needed. \\
Before displaying a custom waveform, it is important to transfer the samples from FLASH to internal RAM, since the read latency of FLASH significantly reduces the output frequency of the waveform.